\documentclass[11pt,a4paper]{article}
\usepackage[T1]{fontenc}
\usepackage{lmodern}
\usepackage{amsmath,amssymb}
\usepackage{booktabs}
\usepackage{array}
\usepackage[a4paper,margin=2.5cm]{geometry}
\usepackage[hidelinks]{hyperref}
\usepackage{amsthm}

\newif\ifshowworked \showworkedtrue
\newenvironment{worked}{\ifshowworked\par\smallskip\noindent\textit{Working.}\par\else\setbox0\vbox\bgroup\fi}{\ifshowworked\par\smallskip\else\egroup\fi}

\newcommand{\Nc}{\mathcal{N}}
\newcommand{\Gc}{\mathcal{G}}
\newcommand{\eps}{\varepsilon}
\newcommand{\ii}{\mathrm{i}}
\newcommand{\bx}{\hat u_x}
\newcommand{\bz}{\hat u_z}
\newcommand{\bw}{\hat w}
\newcommand{\bp}{\hat\varphi}
\newcommand{\bt}{\hat\theta}
\newcommand{\kB}{k_{\mathrm{th}}}
\newcommand{\bwth}{\beta^{w}}

\title{\bfseries Lamb waves in a one-dimensional hexagonal piezoelectric\\[2pt]
quasicrystal plate with Moore--Gibson--Thompson heat conduction\\[2pt]
and the nonlocal strain gradient size effect}
\author{Dr.\ Vipin\\ Gurugram University}
\date{30 October 2026}

\begin{document}
\maketitle

\begin{abstract}
\noindent
A one-dimensional hexagonal quasicrystal has two displacements: the phonon displacement, which is the
ordinary elastic motion, and the phason displacement, which describes the rearrangement of the atomic
clusters and exists only along the quasiperiodic axis. We take that axis along the plate normal $z$, keep
the electric potential and the temperature, and reduce the five coupled field equations to a plane-wave
problem across the thickness. The plate has thickness $2h$ and free faces at $z=\pm h$; the wave is
$\exp[\ii(kx-\omega t)]$ with complex $k$. Heat conduction is of Moore--Gibson--Thompson type and carries a
thermo-phason term; the phonon stress carries the nonlocal strain gradient operator. We give the
constitutive components, the five equations, the $5\times5$ matrix, the degree and parity of its
determinant, the amplitude ratios, the seven conditions per face, the boundary matrix and the hand limits.
\end{abstract}

\noindent\textit{Keywords.} quasicrystal; phason; Lamb wave; Moore--Gibson--Thompson; nonlocal strain
gradient; plate.

\section{Introduction}

A quasicrystal has order without periodicity. In the one-dimensional hexagonal case the atoms sit in planes
that repeat along $x$ and $y$ in the usual way, but along the axis $z$ the spacing follows a sequence that
is ordered and never repeats. Two displacements are needed to describe the motion. The first is the phonon
displacement $u_i$, the same field an ordinary crystal has. The second is the phason displacement $w_z$,
which does not move an atom from one place to another; it describes the local rearrangement of the atomic
clusters along the quasiperiodic direction, the way a queue of people can shuffle along without anyone
leaving the queue. Because the rearrangement costs energy, the phason field carries its own stress, its own
moduli, and its own dynamics.

The dynamics of the phason is the interesting part. It has inertia, so it can oscillate, and it has
friction against the lattice, so it can relax. The two together give a telegraph equation: at high
frequency the phason behaves like a wave, at low frequency like a diffusion. This is written in
Agiasofitou and Lazar \cite[Eqs.~(26) and (30)]{Agiasofitou2014}, who give the friction term in the phason
balance and the telegraph equation that follows from it, and in Baggioli and Landry
\cite[Eq.~(17)]{Baggioli2020}, who write the crossover between the wave and diffusive regimes with a single
relaxation rate. The mechanism behind the friction was proposed by Lubensky, Ramaswamy and Toner
\cite{Lubensky1985}; we cite them for the mechanism only and we did not read their text.

A Lamb wave in such a plate therefore has three partners to talk to: the electric field, the phason
rearrangement and the temperature. Each of them has its own relaxation, and the wave is attenuated when one
of those relaxations falls in step with the wave frequency. This is the reason the problem is worth doing:
the attenuation spectrum carries three separate fingerprints, and a measurement of the attenuation against
frequency can in principle separate them.

Thermoelastic quasicrystal plates have been studied. Zhang, Yu, Zhang and Elmaimouni
\cite{Zhang2021} treat wave propagation in quasicrystal plates and give the constitutive structure we use;
we read their displayed equations as published but did not transcribe their moduli table. Their companion
work on nonlocal quasicrystal nanoplates \cite{Wang2025} and on couple-stress plates \cite{Wang2022} treats
size effects of other kinds. The closest competitor we know of is a 2026 paper on generalised thermoelastic
Lamb waves in one-dimensional hexagonal quasicrystal nanoplates
\cite{ZhangSong2026}. We could not open it; the citation service gave us the title and the authors only. It
is cited here with a note, and the first task of the plan is to obtain it and to write down four things:
which generalised thermoelastic model it uses, whether it has one phase lag or two, whether it reports
attenuation, and whether the phason field is dynamic or static. If their model already contains the
Moore--Gibson--Thompson operator with a thermo-phason term, our thermal novelty is gone and the paper must
lean on the phason dynamics and the size effect instead.

The novelty we claim is the following, in the words agreed with the supervisor. Four choices define it.
The phason field is kept dynamic, with both inertia and relaxation, where most plate work keeps it static.
Heat conduction is of Moore--Gibson--Thompson type and carries the thermo-phason coupling, which to our
knowledge has not been written in a plate problem. Nonlocal strain gradient theory is applied to the phonon
stress, with the phason part kept as a labelled option, and the boundary conditions include the double
tractions of the gradient theory together with a phason traction condition. The parity analysis of
Section~\ref{sec:D} is part of the claim too: it says which couplings decide whether the plate modes split
into symmetric and antisymmetric families, and the answer is not the one we expected.

\section{Setting and assumptions}

The plate occupies $-h\le z\le h$ and is unbounded in $x$ and $y$. Its faces are free of traction, free of
phason traction, free of surface charge and thermally insulating. The material is a one-dimensional
hexagonal piezoelectric quasicrystal. The six-fold axis and the quasiperiodic direction coincide and lie
along $z$.

We take plane strain in the $x$--$z$ plane, so nothing depends on $y$ and $u_y=0$. The wave is
$\exp[\ii(kx-\omega t)]$ with real $\omega$ and complex $k=k_r+\ii k_i$; positive $k_i$ means decay along
$+x$. Every field through the thickness is a sum of $\exp(\eta z)$ with complex $\eta$.

The theory is linear and isothermal in the reference state. Small compared with the lattice parameter is
the phason displacement $w_z(x,z,t)$, so only the phason strains $w_{zz}=\partial_zw_z$ and
$w_{zx}=\partial_xw_z$ are kept. For the electric field we write a potential, $E_i=-\varphi_{,i}$, which is
the quasi-static form. Temperature enters as the change $\theta$ from the reference value $T_0$, small
compared with $T_0$ itself.

Two operators carry the size effect, exactly as in the piezoelectric semiconductor paper:
$\Nc(\eta)=1-l_{\rm nl}^2(\eta^2-k^2)$ on the inertia and $\Gc(\eta)=1-l_g^2(\eta^2-k^2)$ on the phonon
elastic stress. The phason stress carries no size operator in the main model. That is a choice, not a
result; we label the alternative in Section~\ref{sec:tests}.

\section{Notation}\label{sec:not}

\begin{center}\footnotesize
\begin{tabular}{p{2.0cm}p{5.6cm}p{1.9cm}p{4.6cm}}
\toprule
symbol & meaning in words & unit & where it enters\\
\midrule
$h$ & half the plate thickness & m & boundary rows\\
$k,\omega,\eta$ & wave number, angular frequency, thickness rate & rad/m, rad/s, 1/m & everywhere\\
$u_x,u_z$ & phonon displacements along and across the plate & m & rows 1, 2\\
$w$ & phason displacement along the quasiperiodic axis $z$ & m & row 3\\
$\varphi$ & electric potential & V & row 4\\
$\theta$ & temperature change & K & row 5\\
$c_{11},c_{12},c_{13},c_{33},c_{44}$ & phonon elastic stiffnesses & Pa & rows 1, 2\\
$K_1,K_2$ & phason stiffnesses, in-plane and along the axis & Pa & row 3\\
$R_1,R_2,R_3$ & phonon--phason coupling constants & Pa & rows 1--3\\
$e_{31},e_{33},e_{15}$ & piezoelectric constants & C/m$^2$ & rows 1, 2, 4\\
$\epsilon_{11},\epsilon_{33}$ & dielectric permittivities & F/m & row 4\\
$\beta_1,\beta_3$ & thermal stress coefficients of the phonons & Pa/K & rows 1, 2, 5\\
$\bwth$ & thermo-phason coefficient, the phason analogue of $\beta$ & Pa/K & rows 3, 5\\
$p_3$ & pyroelectric constant & C/(m$^2$K) & rows 4, 5\\
$\rho$ & phonon mass density & kg/m$^3$ & rows 1, 2\\
$\rho_p$ & phason mass density & kg/m$^3$ & row 3\\
$\gamma$ & phason friction coefficient, the inverse of a relaxation time & kg/(m$^3$s) & row 3\\
$T_0$ & reference temperature & K & row 5\\
$c_p$, $\kB$ & specific heat, thermal conductivity & J/(kg\,K), W/(m\,K) & row 5\\
$\tau_q$ & MGT relaxation time of the heat flux & s & row 5\\
$l_g$, $l_{\rm nl}$ & gradient and nonlocal lengths of NSGT & m & $\Gc$, $\Nc$\\
$\mathbf M(\eta)$, $\mathbf B$ & amplitude matrix, boundary matrix & --- & Sections~\ref{sec:C}, \ref{sec:G}\\
\bottomrule
\end{tabular}
\end{center}

The everyday picture for the phason is the queue again. Push the queue from behind and it moves as a whole;
that is the phonon. Ask the people in it to close up the gaps without moving the ends; that is the phason.
A real compression costs far more energy in a quasicrystal than this rearrangement does, which is why
$K_1$ and $K_2$ are smaller than the $c$'s, and why the phason branch sits below the phonon branches in
frequency.

\section{Source equations}\label{sec:src}

\subsection{Strain, phason strain, field}

The phonon strain is $\eps_{ij}=\tfrac12(u_{i,j}+u_{j,i})$ and the phason strain is
$w_{ij}=\partial_iw_j$ with only $w_3=w$ non-zero. The electric field is $E_i=-\varphi_{,i}$. Under the plane
wave these give $\eps_{xx}=\ii k\bx$, $\eps_{zz}=\eta\bz$, $\eps_{xz}=\tfrac12(\eta\bx+\ii k\bz)$,
$w_{zz}=\eta\bw$, $w_{zx}=\ii k\bw$, $E_x=-\ii k\bp$, $E_z=-\eta\bp$.
Tag: textbook laws of kinematics and electrostatics; no equation number claimed.

\subsection{Constitutive relations of the quasicrystal}

For the class in question, the phonon stress, the phason stress and the electric displacement are
\begin{equation}\label{eq:const}
\begin{aligned}
\sigma_{ij}&=c_{ijkl}\bigl(1-l_g^2\nabla^2\bigr)\eps_{kl}+R_{ijkl}w_{kl}-e_{kij}E_k-\beta_{ij}\theta,\\
H_{ij}&=R_{klij}\eps_{kl}+K_{ijkl}w_{kl}-\bwth_{ij}\theta,\\
D_i&=e_{ikl}\eps_{kl}+\epsilon_{ik}E_k+p_i\theta .
\end{aligned}
\end{equation}
The first line is the standard phonon constitutive law of a one-dimensional hexagonal quasicrystal, as
written for plates by Zhang and co-workers \cite{Zhang2021}; we read their displayed equations as published
and did not transcribe their table of moduli, so no equation number is quoted from them. Line two is
the phason constitutive law of the same source, again without an equation number from us. Ordinary
piezoelectricity with a pyroelectric term is line three. In the first line, the gradient operator on the
phonon strain is the NSGT placement, and the thermo-phason term $-\bwth_{ij}\theta$ in the second line is ours: it is the phason
analogue of thermal stress and follows from Duhamel's argument applied to the phason strain.
Tag: taken from the source for the elastic, phason and piezoelectric parts; the size operator and the
thermo-phason term are ours.

We must record what we did not read. We did not read a numbered equation for the phason constitutive law in
any source; the structure above is the one used across the quasicrystal literature and is quoted here
without a number. Please locate it in Section~2 of \cite{Zhang2021} and write the number in your own draft.

\subsection{Balance of momentum}

\begin{equation}\label{eq:mom}
\bigl(1-l_{\rm nl}^2\nabla^2\bigr)\rho\,\ddot u_i=\sigma_{ij,j}.
\end{equation}
The differential form is the one quoted in \cite[Eqs.~(18)--(19)]{Barretta2018}, who attribute it to Lim,
Zhang and Reddy \cite{Lim2015}. We did not open \cite{Lim2015}; its equations (17) to (19) are known to us
only through \cite{Barretta2018}, and Section~2 of the original is where you should check it.
Tag: taken from \cite{Barretta2018}; the placement of the operator on the inertia alone is ours.

\subsection{Phason balance}

\begin{equation}\label{eq:phason}
H_{ij,j}=\rho_p\,\ddot w_i+\gamma\,\dot w_i .
\end{equation}
The friction term is taken from \cite[Eq.~(26)]{Agiasofitou2014}, where a damping tensor multiplies the
phason velocity, and the resulting telegraph equation is their Eq.~(30). We keep a scalar $\gamma$ in place
of their tensor, because our geometry has one phason component only; that is our simplification. The
relaxation rate $\gamma/\rho_p$ is written as $1/\tau$ in \cite[Eq.~(17)]{Baggioli2020}.
Tag: taken from the source with the change named, the tensor replaced by a scalar.

\subsection{Gauss's law}

\begin{equation}\label{eq:gauss}
D_{i,i}=0 .
\end{equation}
There are no free carriers in a quasicrystal, so the right side is zero. This is textbook electrostatics of
a dielectric; no equation number is claimed.
Tag: textbook law.

\subsection{MGT heat equation}

\begin{equation}\label{eq:heat}
\Bigl(1+\tau_q\frac{\partial}{\partial t}\Bigr)
\Bigl[\rho c_p\dot\theta+T_0\beta_{ij}\dot\eps_{ij}+T_0\bwth_{ij}\dot w_{ij}
+T_0p_i\dot E_i\Bigr]=\kB\nabla^2\theta .
\end{equation}
The operator is the Moore--Gibson--Thompson operator with one time derivative cancelled, as read in
\cite[Eqs.~(2.5) and (3.1)--(3.5)]{Quintanilla2019}. Classical here are the phonon thermal stress term and the pyroelectric term. The thermo-phason term $T_0\bwth_{ij}\dot w_{ij}$ is ours: it is the reciprocity partner of
the $-\bwth\theta$ term in the phason stress, and it must be there if the entropy production is to be
positive.
Tag: the operator is taken from \cite{Quintanilla2019}; the thermo-phason term is ours, from the entropy
balance.

\subsection{Where the friction term comes from}

The phason balance \eqref{eq:phason} deserves a derivation, because the friction term is the one part of the
model that has no counterpart in ordinary elasticity. Start from the two statements that define the
telegraph behaviour.

\begin{worked}
The first statement is that the phason stress relaxes towards its elastic value with a time constant
$\tau_p=\rho_p/\gamma$; write the relaxation as
$\dot H_{ij}+\frac{1}{\tau_p}\bigl(H_{ij}-K_{ijkl}w_{kl}\bigr)=0$ in the purely phason part. The second
statement is the balance of phason momentum. Combining the two by multiplying the balance by
$(1+\tau_p\partial_t)$ and using the relaxation to replace $\dot H$ gives
\[
K_{ijkl}w_{kl,j}=\rho_p\ddot w_i+\gamma\dot w_i ,
\]
which is the wave--telegraph equation of Agiasofitou and Lazar \cite[Eq.~(30)]{Agiasofitou2014}: one term with
the second time derivative and one with the first. At high frequency the first term dominates and the phason
propagates; at low frequency the second dominates and the phason diffuses. If $\gamma=0$ the equation is a
pure wave equation and the phason never relaxes, which is unphysical for a real quasicrystal.
\end{worked}

\noindent Written for one phason component along $z$ and with the friction collected in the scalar $\gamma$,
this is exactly \eqref{eq:phason} with $K_{ijkl}w_{kl}$ written out as $K_1\partial_x^2w+K_2\partial_z^2w$. The
two constants have a simple reading: $K_1$ is the stiffness against a rearrangement that varies along the
plate, $K_2$ the stiffness against one that varies across the thickness. The mass $\rho_p$ is not a real mass
density in the sense of atoms moving; it is the inertia associated with the collective rearrangement, and its
value is a property of the phason mode, not of the crystal alone.

\subsection{Where the thermo-phason term comes from}

\begin{worked}
Apply Duhamel's argument to the phason strain. Heat the crystal by $\theta$ and the phason strain that would
appear in a free sample is $\alpha^{w}\theta$, with $\alpha^{w}$ a thermo-phason expansion coefficient. If
that strain is prevented, a phason stress $-\bwth\theta$ appears with $\bwth=K\alpha^{w}$, which is the last
term of \eqref{eq:Hzz}. Reciprocity then forces a term $+T_0\bwth\dot w_{zz}$ in the heat equation, because the
same entropy production contains both. That term is the last of the bracket in \eqref{eq:heat}, and it is the
one that produces the second damping peak of Section~\ref{sec:H}.
\end{worked}

\section{Calculation A: constitutive components}\label{sec:A}

\subsection{Which constants survive}

\begin{worked}
The point group is $6mm$ with the axis along $z$, as in the piezoelectric semiconductor paper. A rotation of
$60^\circ$ about $z$ forces $c_{66}=(c_{11}-c_{12})/2$ and $c_{44}=c_{55}$; a rotation of $180^\circ$ about
an in-plane axis is not a symmetry of the poled crystal, so constants with one $z$ index survive. The
quasiperiodic direction is $z$, so the phason strain has only the two components $w_{zz}$ and $w_{zx}$; the
phason moduli therefore reduce to two, $K_{zzzz}=K_2$ and $K_{zxzx}=K_1$, and the phonon--phason coupling to
three, $R_{xxzz}=R_1$, $R_{zzzz}=R_2$ and $R_{xzxz}=R_3$. In plane strain with $\eps_{yy}=0$ the constant
$c_{12}$ drops out of every equation, exactly as in the semiconductor paper.
\end{worked}

\noindent The constants that appear are $c_{11},c_{13},c_{33},c_{44}$; $K_1,K_2$; $R_1,R_2,R_3$;
$e_{31},e_{33},e_{15}$; $\epsilon_{11},\epsilon_{33}$; $\beta_1,\beta_3$; $\bwth$; and $p_3$. The two
piezoelectric-phason couplings that some papers write are absent in our model, because the phason
rearrangement does not separate charge in the class we treat; if your source keeps such a coupling, say so
and add one entry to row 4.

\subsection{Stresses and induction in components}

\begin{worked}
Apply the same routine as in the semiconductor paper. Each $\partial_x$ becomes $\ii k$, each $\partial_z$
becomes $\eta$, and the Laplacian becomes $\eta^2-k^2$. For $\sigma_{xx}$: $c_{11}\eps_{xx}$ gives
$c_{11}\ii k\bx$; $c_{13}\eps_{zz}$ gives $c_{13}\eta\bz$; the phason coupling $R_1w_{zz}$ gives $R_1\eta\bw$;
the piezoelectric term $-e_{31}E_z$ gives $+e_{31}\eta\bp$; the thermal term is $-\beta_1\bt$. The gradient
operator multiplies the two phonon elastic terms only. For $H_{zz}$: $R_1\eps_{xx}$ gives $R_1\ii k\bx$,
$R_2\eps_{zz}$ gives $R_2\eta\bz$, $K_2w_{zz}$ gives $K_2\eta\bw$, and the thermo-phason term gives
$-\bwth\bt$. For $D_z$: $e_{31}\eps_{xx}$ gives $e_{31}\ii k\bx$, $e_{33}\eps_{zz}$ gives $e_{33}\eta\bz$,
$-\epsilon_{33}E_z$ gives $+\epsilon_{33}\eta\bp$, and the pyroelectric term gives $+p_3\bt$.
\end{worked}

\noindent The seven constitutive components are
\begin{align}
\hat\sigma_{xx}&=\Gc\bigl(\ii kc_{11}\bx+\eta c_{13}\bz\bigr)+R_1\eta\bw+e_{31}\eta\bp-\beta_1\bt,
\label{eq:sxx}\\
\hat\sigma_{zz}&=\Gc\bigl(\ii kc_{13}\bx+\eta c_{33}\bz\bigr)+R_2\eta\bw+e_{33}\eta\bp-\beta_3\bt,
\label{eq:szz}\\
\hat\sigma_{xz}&=\Gc c_{44}\bigl(\eta\bx+\ii k\bz\bigr)+R_3\ii k\bw+\ii ke_{15}\bp,
\label{eq:sxz}\\
\hat H_{zz}&=R_1\ii k\bx+R_2\eta\bz+K_2\eta\bw-\bwth\bt,
\label{eq:Hzz}\\
\hat H_{zx}&=R_3\bigl(\eta\bx+\ii k\bz\bigr)+K_1\ii k\bw,
\label{eq:Hzx}\\
\hat D_x&=e_{15}\bigl(\eta\bx+\ii k\bz\bigr)-\ii k\epsilon_{11}\bp,
\label{eq:Dx}\\
\hat D_z&=\ii ke_{31}\bx+\eta e_{33}\bz-\eta\epsilon_{33}\bp+p_3\bt .
\label{eq:Dz}
\end{align}

Notice where the phason sits. It couples to the phonon through $R_1$ and $R_2$ in the normal stresses and
through $R_3$ in the shear, and its own stress is driven by the phonon strain through the same constants.
That is the reciprocity of the phonon--phason coupling, and it is why the matrix of Section~\ref{sec:C} is
not triangular: the phason equation feels the phonon and the phonon equations feel the phason.

\section{Calculation B: the five governing equations}\label{sec:B}

\subsection{Momentum}

\begin{worked}
The $x$ component of \eqref{eq:mom} is $-\bigl(\ii k\hat\sigma_{xx}+\eta\hat\sigma_{xz}\bigr)-\Nc\rho\omega^2\bx=0$,
because the divergence is $\partial_x\sigma_{xx}+\partial_z\sigma_{xz}$ and $\ddot u_x=-\omega^2\bx$. The $z$
component is $-\bigl(\ii k\hat\sigma_{xz}+\eta\hat\sigma_{zz}\bigr)-\Nc\rho\omega^2\bz=0$.
\end{worked}

\noindent
\begin{equation}\label{eq:g1}
-\bigl(\ii k\hat\sigma_{xx}+\eta\hat\sigma_{xz}\bigr)-\Nc\rho\omega^2\bx=0,\qquad
-\bigl(\ii k\hat\sigma_{xz}+\eta\hat\sigma_{zz}\bigr)-\Nc\rho\omega^2\bz=0 .
\end{equation}

\subsection{Phason balance}

\begin{worked}
The phason balance \eqref{eq:phason} has one free index, $i=3$. Its left side is
$H_{3j,j}=\partial_xH_{zx}+\partial_zH_{zz}$, which under the plane wave is $\ii k\hat H_{zx}+\eta\hat H_{zz}$.
The right side is $\rho_p\ddot w+\gamma\dot w=-\rho_p\omega^2\bw-\ii\omega\gamma\bw$. Bringing everything to
one side gives $-\bigl(\ii k\hat H_{zx}+\eta\hat H_{zz}\bigr)-\rho_p\omega^2\bw-\ii\omega\gamma\bw=0$, and
changing the overall sign for convenience,
\[
\bigl(\ii k\hat H_{zx}+\eta\hat H_{zz}\bigr)+\rho_p\omega^2\bw+\ii\omega\gamma\bw=0 .
\]
\end{worked}

\noindent
\begin{equation}\label{eq:g2}
\bigl(\ii k\hat H_{zx}+\eta\hat H_{zz}\bigr)+\bigl(\rho_p\omega^2+\ii\omega\gamma\bigr)\bw=0 .
\end{equation}

\subsection{Gauss's law}

\begin{worked}
With \eqref{eq:gauss}, $\ii k\hat D_x+\eta\hat D_z=0$.
\end{worked}

\noindent
\begin{equation}\label{eq:g3}
\ii k\hat D_x+\eta\hat D_z=0 .
\end{equation}

\subsection{Heat balance}

\begin{worked}
Put the plane wave into \eqref{eq:heat}. The time derivative is $-\ii\omega$, so the bracket is
$-\ii\omega\rho c_p\bt+T_0\beta_1(-\ii\omega)(\ii k\bx)+T_0\beta_3(-\ii\omega)(\eta\bz)
+T_0\bwth(-\ii\omega)(\eta\bw)+T_0p_3(-\ii\omega)(-\eta\bp)$, which equals
$-\ii\omega\rho c_p\bt+T_0\omega k\beta_1\bx-\ii\omega T_0\beta_3\eta\bz-\ii\omega T_0\bwth\eta\bw
+\ii\omega T_0p_3\eta\bp$. Multiply by $f=1-\ii\omega\tau_q$ and subtract $\kB(\eta^2-k^2)\bt$.
\end{worked}

\noindent
\begin{equation}\label{eq:g4}
f\Bigl[-\ii\omega\rho c_p\bt+T_0\omega k\beta_1\bx-\ii\omega T_0\beta_3\eta\bz
-\ii\omega T_0\bwth\eta\bw+\ii\omega T_0p_3\eta\bp\Bigr]-\kB L\bt=0,\qquad L=\eta^2-k^2 .
\end{equation}

\section{Calculation C: the amplitude matrix}\label{sec:C}

The unknowns are ordered as $\mathbf a=(\bx,\bz,\bw,\bp,\bt)^{\top}$. Rows are, in order, momentum $x$,
momentum $z$, phason, Gauss, heat.

\subsection{Row 1, momentum along the plate}

\begin{worked}
Put \eqref{eq:sxx} and \eqref{eq:sxz} into the first of \eqref{eq:g1}. For $\bx$: $\ii k\cdot\Gc c_{11}\ii k=-k^2\Gc c_{11}$
and $\eta\cdot\Gc c_{44}\eta=\eta^2\Gc c_{44}$; with the overall minus and the inertia,
$M_{11}=-\Nc\rho\omega^2+\Gc(c_{11}k^2-c_{44}\eta^2)$. For $\bz$: $\ii k\Gc c_{13}\eta+\eta\Gc c_{44}\ii k=\ii k\eta\Gc(c_{13}+c_{44})$,
with a minus. For $\bw$: $\ii kR_1\eta+\eta R_3\ii k=\ii k\eta(R_1+R_3)$, with a minus. For $\bp$:
$\ii k e_{31}\eta+\eta\ii k e_{15}=\ii k\eta(e_{31}+e_{15})$, with a minus. For $\bt$: $\ii k(-\beta_1)$, with
the overall minus, gives $+\ii k\beta_1$.
\end{worked}

\noindent
\begin{equation}\label{eq:row1}
M_{11}=-\Nc\rho\omega^2+\Gc(c_{11}k^2-c_{44}\eta^2),\quad
M_{12}=-\ii k\eta\Gc(c_{13}+c_{44}),\quad
M_{13}=-\ii k\eta(R_1+R_3),
\end{equation}
\begin{equation}\label{eq:row1b}
M_{14}=-\ii k\eta(e_{31}+e_{15}),\qquad M_{15}=\ii k\beta_1 .
\end{equation}

\subsection{Row 2, momentum across the plate}

\begin{worked}
For $\bx$: $\ii k\Gc c_{44}\eta+\eta\Gc c_{13}\ii k=\ii k\eta\Gc(c_{13}+c_{44})$, with a minus. For $\bz$:
$\ii k\Gc c_{44}\ii k=-k^2\Gc c_{44}$ and $\eta\Gc c_{33}\eta=\eta^2\Gc c_{33}$, so
$M_{22}=-\Nc\rho\omega^2+\Gc(c_{44}k^2-c_{33}\eta^2)$. For $\bw$: $\ii kR_3\ii k=-k^2R_3$ and $\eta R_2\eta=\eta^2R_2$;
with the overall minus, $M_{23}=R_3k^2-R_2\eta^2$. For $\bp$: $\ii k\ii ke_{15}=-k^2e_{15}$ and $\eta\eta e_{33}=\eta^2e_{33}$,
with a minus. For $\bt$: $\eta(-\beta_3)$ with the minus gives $+\eta\beta_3$.
\end{worked}

\noindent
\begin{equation}\label{eq:row2}
M_{21}=-\ii k\eta\Gc(c_{13}+c_{44}),\quad
M_{22}=-\Nc\rho\omega^2+\Gc(c_{44}k^2-c_{33}\eta^2),\quad
M_{23}=R_3k^2-R_2\eta^2,
\end{equation}
\begin{equation}\label{eq:row2b}
M_{24}=e_{15}k^2-e_{33}\eta^2,\qquad M_{25}=\eta\beta_3 .
\end{equation}

\subsection{Row 3, the phason equation}

\begin{worked}
Use \eqref{eq:Hzz} and \eqref{eq:Hzx} in \eqref{eq:g2}, which has no overall minus, so the entries keep the
signs of the divergence. For $\bx$: $\ii kR_3\eta+\eta R_1\ii k=\ii k\eta(R_1+R_3)$. For $\bz$:
$\ii kR_3\ii k=-k^2R_3$ and $\eta R_2\eta=\eta^2R_2$, so $M_{32}=R_2\eta^2-R_3k^2$.
For $\bw$: $\ii kK_1\ii k=-k^2K_1$ and $\eta K_2\eta=\eta^2K_2$, plus the inertia and friction
$\rho_p\omega^2+\ii\omega\gamma$ from \eqref{eq:g2}; so $M_{33}=-K_1k^2+K_2\eta^2+\rho_p\omega^2+\ii\omega\gamma$.
For $\bp$: nothing, the phason does not see the electric field in our model. For $\bt$: $\eta(-\bwth)=-\eta\bwth$.
\end{worked}

\noindent
\begin{equation}\label{eq:row3}
M_{31}=\ii k\eta(R_1+R_3),\quad
M_{32}=R_2\eta^2-R_3k^2,\quad
M_{33}=K_2\eta^2-K_1k^2+\rho_p\omega^2+\ii\omega\gamma,
\end{equation}
\begin{equation}\label{eq:row3b}
M_{34}=0,\qquad M_{35}=-\eta\bwth .
\end{equation}

\subsection{Rows 4 and 5}

\begin{worked}
Row 4 is Gauss: $\ii k\hat D_x+\eta\hat D_z$ with \eqref{eq:Dx} and \eqref{eq:Dz}. For $\bx$:
$\ii ke_{15}\eta+\eta\ii ke_{31}=\ii k\eta(e_{15}+e_{31})$. For $\bz$: $-k^2e_{15}+\eta^2e_{33}$. For $\bw$:
zero. For $\bp$: $\ii k(-\ii k\epsilon_{11})+\eta(-\eta\epsilon_{33})=\epsilon_{11}k^2-\epsilon_{33}\eta^2$.
For $\bt$: $\eta p_3$. Row 5 is \eqref{eq:g4} read off directly.
\end{worked}

\noindent
\begin{equation}\label{eq:row4}
M_{41}=\ii k\eta(e_{15}+e_{31}),\quad M_{42}=e_{33}\eta^2-e_{15}k^2,\quad M_{43}=0,\quad
M_{44}=\epsilon_{11}k^2-\epsilon_{33}\eta^2,\quad M_{45}=\eta p_3,
\end{equation}
\begin{equation}\label{eq:row5}
M_{51}=fT_0\omega k\beta_1,\quad M_{52}=-\ii\omega fT_0\beta_3\eta,\quad
M_{53}=-\ii\omega fT_0\bwth\eta,\quad M_{54}=\ii\omega fT_0p_3\eta,\quad
M_{55}=-\ii\omega f\rho c_p-\kB L .
\end{equation}

\subsection{The matrix}

Collecting \eqref{eq:row1}--\eqref{eq:row5},
\begin{equation}\label{eq:M}
\mathbf M(\eta)\mathbf a=\mathbf 0,\qquad
\mathcal N=1-l_{\rm nl}^2L,\quad \mathcal G=1-l_g^2L,\quad L=\eta^2-k^2,\quad f=1-\ii\omega\tau_q .
\end{equation}
The phason enters in exactly five places: $M_{13}$, $M_{23}$, $M_{31}$, $M_{32}$ through the coupling
constants, and $M_{33}$ through its own stiffness, inertia and friction. The thermo-phason coefficient sits in
$M_{35}$ and $M_{53}$, one each side, which is the reciprocity of the pair.

\section{Calculation D: degree and parity}\label{sec:D}

\subsection{Degree}

\begin{worked}
Write the highest power of $\eta$ in each entry. Row 1: $M_{11}\sim\eta^4$ when $l_g\ne0$ (from
$\Gc\eta^2\cdot c_{44}\eta^2$), $M_{12}\sim\eta^3$, $M_{13}\sim\eta^1$, $M_{14}\sim\eta^1$, $M_{15}\sim\eta^0$.
Row 2: $M_{21}\sim\eta^3$, $M_{22}\sim\eta^4$, $M_{23}\sim\eta^2$, $M_{24}\sim\eta^2$, $M_{25}\sim\eta^1$.
Row 3: $M_{31}\sim\eta^1$, $M_{32}\sim\eta^2$, $M_{33}\sim\eta^2$, $M_{35}\sim\eta^1$.
Row 4: $M_{41}\sim\eta^1$, $M_{42}\sim\eta^2$, $M_{44}\sim\eta^2$, $M_{45}\sim\eta^1$.
Row 5: $M_{51}\sim\eta^0$, $M_{52}\sim\eta^1$, $M_{53}\sim\eta^1$, $M_{54}\sim\eta^1$, $M_{55}\sim\eta^2$.
The permutation with the largest total takes $M_{11}$ (4), $M_{22}$ (4), $M_{33}$ (2), $M_{44}$ (2) and
$M_{55}$ (2), giving 14. No other permutation exceeds it, because rows 1 and 2 contribute at most four each
and rows 3, 4 and 5 at most two each.
\end{worked}

\noindent
\begin{equation}\label{eq:deg}
\deg_\eta\det\mathbf M=\begin{cases}14,&l_g\ne0,\\10,&l_g=0 .\end{cases}
\end{equation}
The phason equation is of second order in $z$, like the others, so it adds no roots and no conditions. The
gradient operator adds four roots, and with them four conditions, two per face.

\subsection{Parity}

\begin{worked}
Again suppose a sign vector $s=(s_1,\dots,s_5)$ with $p(i,j)=s_is_j$, where $p$ is the parity of an entry.
From row 1: $M_{11}$ even gives nothing; $M_{12}$ odd gives $s_2=-s_1$; $M_{13}$ odd gives $s_3=-s_1$;
$M_{14}$ odd gives $s_4=-s_1$; $M_{15}$ even gives $s_5=s_1$. Take $s_1=-1$. Then
$s=(-1,+1,+1,+1,-1)$. Check the rest. Row 2: $M_{21}$ odd needs $s_2s_1=-1$, and $+1\cdot-1=-1$, fine;
$M_{23}$ even needs $s_2s_3=+1$, and it is, fine; $M_{24}$ even, fine; $M_{25}$ odd needs $s_2s_5=-1$, fine.
Row 3: $M_{31}$ odd needs $s_3s_1=-1$, fine; $M_{32}$ even, fine; $M_{33}$ even, fine; $M_{35}$ odd needs
$s_3s_5=-1$, fine. Row 4: $M_{41}$ odd needs $s_4s_1=-1$, fine; $M_{42}$ even, fine; $M_{44}$ even, fine;
$M_{45}$ odd needs $s_4s_5=-1$, fine. Row 5: $M_{51}$ even needs $s_5s_1=+1$, fine; $M_{52}$ odd,
$M_{53}$ odd, $M_{54}$ odd all need products $-1$, fine; $M_{55}$ even, fine.
No entry fails. So the sign vector always exists.
\end{worked}

\noindent
\begin{equation}\label{eq:parity}
\det\mathbf M(-\eta)=\det\mathbf M(\eta)\quad\text{for every value of every coupling constant,}
\end{equation}
with the sign vector $(-1,+1,+1,+1,-1)$ in the order $(\bx,\bz,\bw,\bp,\bt)$: the phonon displacement along
the plate is odd through the thickness, and so is the temperature; the other three fields are even.

This is the result we did not expect, and you should test it yourself before you write it. The phason
coupling constants $R_1,R_2,R_3$ are not small, and one would guess that a field that couples the normal
stress to a rearrangement along the axis must break the mirror symmetry. It does not. The reason is that
$R_1$ and $R_2$ always appear as an odd number of $z$ derivatives in the momentum equations and an even
number in the phason equation, which is exactly the pattern the sign vector needs. The thermo-phason term
does the same: it is even in the phason equation and odd in the heat equation, again matching. What breaks
the pairing in the semiconductor paper is not a polar coupling but the carrier field: a density that is even
sits in a row coupled to a temperature that must be odd, and no sign vector can satisfy both. Remove the
carrier field and the obstruction disappears.

The consequence is practical. Plate modes of a quasicrystal plate split into a symmetric and an
antisymmetric family, and the split survives every coupling we have. Its algebraic form is a factorisation
of the boundary matrix, $\det\mathbf B=\det\mathbf B_{\rm S}\det\mathbf B_{\rm A}$, so the two families
can be found separately. That is a strong simplification, and it also means that any observed mixing of the two families
in an experiment cannot come from the couplings in this model; it must come from the geometry, from a
lateral boundary, or from a nonlinearity.

\section{Calculation E: amplitude ratios and partial waves}\label{sec:E}

\begin{worked}
For a root $\eta_m$ the matrix $\mathbf M(\eta_m)$ is singular. The adjugate $\mathbf A$ of
$\mathbf M$ satisfies $\mathbf M\mathbf A=\det(\mathbf M)\mathbf I$, so every column of $\mathbf A$ vanishes
when multiplied by $\mathbf M$ at the root. Compute the entries of the fifth row of the adjugate as
$A_{5j}=(-1)^{5+j}$ times the $4\times4$ minor obtained by deleting row $j$ and column $5$, take the column
with the largest entries, and divide by its largest entry. The four ratios $\rho^{(1)}_m,\dots,\rho^{(4)}_m$
are then finite, and the fifth entry of the column is one by construction.
\end{worked}

\noindent The general solution is
\begin{equation}\label{eq:general}
\bigl(u_x,u_z,w,\varphi,\theta\bigr)=\sum_{m=1}^{N}a_m
\bigl(\rho^{(1)}_m,\rho^{(2)}_m,\rho^{(3)}_m,\rho^{(4)}_m,1\bigr)e^{\eta_mz}e^{\ii(kx-\omega t)},
\end{equation}
with $N=14$ when $l_g\ne0$ and $N=10$ when $l_g=0$. The number of free constants $a_m$ equals the number of
boundary conditions we may impose, seven per face when $l_g\ne0$ and five per face when $l_g=0$.

One consequence of the parity \eqref{eq:parity} deserves a line here. When the sign vector exists, the
roots come in pairs $\pm\eta_m$ and the amplitude ratios of the pair are related: the odd fields change sign
and the even fields keep it. So each partial wave has a twin, and the pair can be combined into a
hyperbolic cosine for the even fields and a hyperbolic sine for the odd ones. That is the mathematical
statement behind the split, and it is worth checking on one pair of roots before you trust the algebra.

\section{Calculation F: boundary conditions}\label{sec:F}

At each face, with outward normal $\mathbf n=(0,0,\pm1)$:

\begin{enumerate}
\item[(b1)] normal traction, with the surface gradient of the double traction:
$\bigl[\sigma_{zz}-\partial_xR_{xz}\bigr]_{z=z_0}=0$;
\item[(b2)] shear traction, with its surface gradient:
$\bigl[\sigma_{xz}-\partial_xR_{xx}\bigr]_{z=z_0}=0$;
\item[(b3)] normal double traction: $\bigl[R_{zz}\bigr]_{z=z_0}=0$;
\item[(b4)] in-plane double traction: $\bigl[R_{zx}\bigr]_{z=z_0}=0$;
\item[(b5)] phason traction: $\bigl[H_{zz}\bigr]_{z=z_0}=0$;
\item[(b6)] open circuit: $\bigl[D_z\bigr]_{z=z_0}=0$;
\item[(b7)] adiabatic face: $\bigl[\partial_z\theta\bigr]_{z=z_0}=0$.
\end{enumerate}

The surface-gradient form of (b1) and (b2) and the double-traction conditions (b3), (b4) are taken from the
gradient elasticity treatment of Charalambopoulos, Douka and Mavratzas
\cite[Sec.~2, Eq.~(14)]{Charalambopoulos2019}, extended by us to the phason, electric and thermal parts.
Conditions (b5), (b6) and (b7) are our modelling choices; the phason condition states that no phason
traction acts on the free face, which is the natural statement for a face where the clusters are free to
rearrange. If your source uses a different phason condition, for example a clamped phason displacement
$w=0$ at the face, say so; the change is one row.

\begin{worked}
The nonlocal operator is inverted root by root, for the same reason as in the semiconductor paper: the
constitutive law at root $\eta_m$ reads $\Nc(\eta_m)\hat\sigma_m=\hat\sigma^{\rm loc}_m$, so the stress that
enters a boundary row is $\hat\sigma^{\rm loc}_m/\Nc(\eta_m)$, and the factor sits inside column $m$. No
common operator can be pulled out of the sum over roots.
\end{worked}

\noindent Evaluating \eqref{eq:sxx}--\eqref{eq:Dz} with the amplitudes of partial wave $m$, dividing the
phonon stresses by $\Nc(\eta_m)$, and using $\partial_x\to\ii k$, the seven row functions at $z=+h$ are
\begin{align}
b_1&=\frac{1}{\Nc(\eta_m)}\Bigl[\hat\sigma_{zz}-\ii k\,l_g^2c_{44}L\hat\eps_{xz}\Bigr],\quad
b_2=\frac{1}{\Nc(\eta_m)}\Bigl[\hat\sigma_{xz}-\ii k\,l_g^2c_{11}L\hat\eps_{xx}\Bigr],
\label{eq:b12}\\
b_3&=\frac{1}{\Nc(\eta_m)}l_g^2c_{33}L\hat\eps_{zz},\quad
b_4=\frac{1}{\Nc(\eta_m)}l_g^2c_{44}L\hat\eps_{xz},\quad
b_5=\hat H_{zz},\quad b_6=\hat D_z,\quad b_7=\eta_m\bt .
\label{eq:b34567}
\end{align}

The phason row $b_5$ is the one to watch. It carries $R_1\ii k\bx+R_2\eta\bz+K_2\eta\bw-\bwth\bt$, and it is
the only row in which the thermo-phason coefficient appears alone; if you want to see the size of the
thermo-phason effect on a dispersion curve, switch this coefficient off and look at how much the phason
traction row changes.

\subsection{The two faces}

\begin{worked}
At $z=-h$ the partial wave carries $\exp(-\eta_mh)$ instead of $\exp(\eta_mh)$. The outward normal is
$(0,0,-1)$, so the normal components flip sign: $\sigma_{zz}$, $\sigma_{xz}$, $H_{zz}$ and $D_z$ all appear
with a minus. The double tractions are quadratic in the normal and keep their sign. The temperature gradient
changes sign because $\partial_z$ is odd.
\end{worked}

\noindent With $c_m=\exp(\eta_mh)$ and $d_m=\exp(-\eta_mh)$ the fourteen rows are
\begin{equation}\label{eq:rows}
b_1c_m,\ b_2c_m,\ b_3c_m,\ b_4c_m,\ b_5c_m,\ b_6c_m,\ b_7c_m\quad(z=+h),\qquad
-b_1d_m,\ -b_2d_m,\ b_3d_m,\ b_4d_m,\ -b_5d_m,\ -b_6d_m,\ -b_7d_m\quad(z=-h).
\end{equation}

\section{Calculation G: the boundary matrix}\label{sec:G}

\begin{worked}
Number the roots $\eta_1,\dots,\eta_N$. Put the row functions of \eqref{eq:b12} and \eqref{eq:b34567} in the
fourteen rows of \eqref{eq:rows} and the roots in fourteen columns. Every entry is a product of a material
constant, an amplitude ratio and an exponential factor. The determinant of the resulting matrix vanishes
exactly when the plate carries a wave.
\end{worked}

\noindent
\begin{equation}\label{eq:disp}
\det\mathbf B(k,\omega)=0,\qquad \mathbf B \text{ is }14\times14\ (l_g\ne0),\ 10\times10\ (l_g=0).
\end{equation}
When $l_g=0$ the rows (b3) and (b4) vanish identically, because each of their entries carries $l_g^2$; they
must be deleted with the four extra roots.

Because the parity \eqref{eq:parity} holds, the matrix blocks. Draw the odd amplitude from each partial wave
into one set and the even amplitude into the other; the rows separate the same way, and
\begin{equation}\label{eq:factor}
\det\mathbf B=\det\mathbf B_{\rm S}\cdot\det\mathbf B_{\rm A},
\end{equation}
with seven rows and seven columns in each block. If you build the matrix and this factorisation does not
appear, either a row has the wrong parity or the roots have not been paired correctly.

The frequency $\omega$ is real and given, the wave number $k$ is complex, and \eqref{eq:disp} is one complex
equation in two real unknowns. Before any determinant is evaluated, multiply column $m$ by
$\exp(-|\mathrm{Re}\,\eta_m|h)$. That factor is common to a column and does not move a zero, but it keeps the
columns within a few orders of magnitude of each other, which they otherwise are not once the gradient roots
appear.

\section{Calculation H: hand limits}\label{sec:H}

\subsection{Phason off: back to the piezoelectric thermoelastic plate}

\begin{worked}
Set $R_1=R_2=R_3=0$, $\bwth=0$ and $\bw=0$. Row 3 and column 3 disappear from the matrix, and
$\det\mathbf M$ becomes the $4\times4$ determinant of the phonon, electric and thermal fields. Put
$e_{31}=e_{33}=e_{15}=0$, $p_3=0$ and $\beta_1=\beta_3=0$ as well, and the $2\times2$ mechanical block
remains:
\[
\begin{vmatrix}-\rho\omega^2+c_{11}k^2-c_{44}\eta^2&-\ii k\eta(c_{13}+c_{44})\\
-\ii k\eta(c_{13}+c_{44})&-\rho\omega^2+c_{44}k^2-c_{33}\eta^2\end{vmatrix}=0 .
\]
With $c_{11}=c_{33}=\lambda+2\mu$ and $c_{13}=\lambda$, $c_{44}=\mu$, this is the classical Rayleigh--Lamb
problem, and the boundary conditions reduce to the two traction conditions per face.
\end{worked}

\noindent The two families then obey the classical relations
\begin{equation}\label{eq:RL}
\frac{\tan(qh)}{\tan(ph)}=-\frac{4k^2pq}{(k^2-q^2)^2}\ ({\rm symmetric}),\qquad
\frac{\tan(qh)}{\tan(ph)}=-\frac{(k^2-q^2)^2}{4k^2pq}\ ({\rm antisymmetric}),
\end{equation}
with $p^2=\omega^2/c_L^2-k^2$ and $q^2=\omega^2/c_T^2-k^2$. This is the check that the phonon part of the
matrix is right, and it is the first test in the list.

\subsection{Phason relaxation: wave or diffusion}

\begin{worked}
Switch off everything except the phason and take the long-wavelength limit, so that $k$ is small. Row 3 reads
$K_2\eta^2+\rho_p\omega^2+\ii\omega\gamma=0$, that is
\[
\eta^2=-\frac{\rho_p\omega^2+\ii\omega\gamma}{K_2}.
\]
At high frequency, $\rho_p\omega\gg\gamma$, the right side is $-\rho_p\omega^2/K_2$ and the phason travels as
a wave with speed $\sqrt{K_2/\rho_p}$. At low frequency, $\gamma\gg\rho_p\omega$, the right side is
$-\ii\omega\gamma/K_2$, which is the dispersion of a diffusion with diffusivity $K_2/\gamma$. The crossover
frequency is $\omega_c=\gamma/\rho_p$, the inverse of the phason relaxation time $\tau_p=\rho_p/\gamma$.
\end{worked}

\noindent
\begin{equation}\label{eq:phasonspeed}
c_{\rm ph}=\sqrt{\frac{K_2}{\rho_p}}\ \ (\omega\tau_p\gg1),\qquad
D_{\rm ph}=\frac{K_2}{\gamma}\ \ (\omega\tau_p\ll1),\qquad \tau_p=\frac{\rho_p}{\gamma}.
\end{equation}
This is the second test: a plot of the phason branch against frequency must flatten to the diffusive form
below $\omega_c$ and become a straight line of slope $c_{\rm ph}$ above it. If your numerical branch does
neither, look at the signs in $M_{33}$.

\subsection{Thermo-phason coupling: the damping partner}

\begin{worked}
Keep only the phason and the temperature. The two rows read
$(-\ii\omega\rho c_p)f\bt-\kB L\bt-\ii\omega fT_0\bwth\eta\bw=0$ and
$(K_2\eta^2+\rho_p\omega^2+\ii\omega\gamma)\bw-\eta\bwth\bt=0$. Eliminating $\bw$ gives a complex correction
to $K_2$ of order $\bwth^2/(\rho c_p)$ times the thermal relaxation factor
$\ii\omega f/(\ii\omega f\rho c_p+\kB L)$, which has the same shape as the thermoelastic damping of an
ordinary solid, with the phason relaxation time playing the role of the thermal one.
\end{worked}

\noindent The thermo-phason term therefore produces a second damping peak, at a frequency set by
$\omega\tau_p\sim1$ rather than by the thermal relaxation. Two peaks in the attenuation spectrum, one from
the temperature and one from the phason, are the observable signature of this model. Whether they can be
separated experimentally depends on their width, and that is a question for the results section, not for
this file.

\subsection{The two size lengths}

\begin{worked}
Set $l_g=l_{\rm nl}=0$ and the classical plate returns exactly, with $\Nc=\Gc=1$ and ten roots. Turning $l_g$
on stiffens the phonon elastic terms, since $\Gc>1$ for a short wave; turning $l_{\rm nl}$ on softens the
inertia, since $\Nc<1$ there. The two shifts are in opposite directions, which is the point of NSGT: a
stiffening and a softening scale, each of them a length, and only their ratio is measured by a dispersion
curve.
\end{worked}

\noindent The phason part carries no size operator in this model, so the phason branch is the cleanest place
to look for the phonon size effect: any shift of the phason branch with $l_g$ or $l_{\rm nl}$ must come
through the coupling $R$, and its size tells you how strongly the two fields are tied.

\section{Self-test list}\label{sec:tests}

\begin{enumerate}
\item Units within every row of \eqref{eq:M}. Row 3 is in Pa/m; check that each of its five entries is.
\item Degree: fourteen with $l_g\ne0$, ten with $l_g=0$.
\item Parity: the sign vector $(-1,+1,+1,+1,-1)$ must work for every choice of the coupling constants. Test
two extreme cases, with all couplings on and with all couplings off.
\item The classical limit \eqref{eq:RL}, both families, and the three long- and short-wavelength limits of
the ordinary plate.
\item The phason speed and diffusivity of \eqref{eq:phasonspeed}, including the crossover.
\item The factorisation \eqref{eq:factor}, checked on the assembled matrix and not on the determinant of the
full matrix alone.
\item The number of roots equal to the number of rows of $\mathbf B$, before and after the switch $l_g=0$.
\item The sign of the thermo-phason effect: switching $\bwth$ on must move the phason branch and the
attenuation, and must do nothing to the purely elastic limit.
\item The alternative boundary conditions: clamped phason displacement, closed circuit, isothermal face, one
at a time, each changing one row.
\item The size effect: stiffening with $l_g$, softening with $l_{\rm nl}$, and a phason branch that shifts
only through the coupling.
\end{enumerate}

\section{Work plan}

\begin{enumerate}
\item Week 1. Obtain the 2026 paper of Zhang, Song, Li, Han, Yu and Yang and write down its model, its heat
equation, its phason treatment and whether it reports attenuation. Until this is done, the novelty paragraph
stays provisional.
\item Week 2. Read Lim, Zhang and Reddy Section 2 and Barretta and Marotti de Sciarra equations (18) and
(19). Confirm the operator placement.
\item Week 3. Read the phason balance in Agiasofitou and Lazar around equations (26) and (30), and the
crossover in Baggioli and Landry equation (17). Write the scalar friction case yourself.
\item Week 4. Redo Calculations A to D without looking at this file, and check the parity statement twice.
\item Week 5. Redo Calculations E to G, build the matrix, and run the ten self-tests.
\item Week 6. Sensitivity: the two sizes, the phason relaxation and the thermo-phason coefficient, one at a
time. Write the paper.
\end{enumerate}

\bibliographystyle{unsrt}
\bibliography{Lamb_Neha}

\end{document}
